%
\againframe{plan}
%
\begin{frame}{Problem at hand}
	%
	\begin{figure}
		\includegraphics[width=0.8\textwidth]{metals_coulomb_force}
	\end{figure}
	\begin{itemize}
		\item The problem at hand is to solve the motion of interacting particles-- electrons of the order of \num{e23} and ions of the order of \num{e23}.
		\item This is a tough problem. The idea is to simplify the problem by making \textbf{approximations}.
	\end{itemize}
	%
\end{frame}
\begin{frame}{Classical free electron theory}
	%
	\begin{columns}
		\begin{column}[t]{0.85\textwidth}
			\begin{itemize}
				\item Ohm's law was discovered in 1827 and needed explanation.
				\item Proposed by Paul Drude in 1900 using analogy of gas of molecules.
				\item Mutual repulsion between electrons is neglected i.e. electrons are independent -- \textbf{independent electron approximation}. 
				\item Assumed ``gas'' is \textbf{free} i.e. not under influence of  lattice --  \textbf{free electron approximation}

				\item Role of lattice is to redistribute the velocity distribution that obeys \textbf{classical} Maxwell-Boltzmann statistics --  classical thermodynamics.
				\item \fcolorbox{green}{white}{The theory explains Ohm's law}
				\item Hendrik Lorentz extended theory to explain thermal conductivity -- hence Drude-Lorentz electron theory.
			\end{itemize}
		\end{column}
		%
		\begin{column}[t]{0.13\textwidth}
			\begin{figure}
				\centering
				\includegraphics[width=0.9\textwidth]{Georg_Simon_Ohm_1789_1854}
				\\
				{\centering\footnotesize G. Ohm}
				\includegraphics[width=0.9\textwidth]{Paul_Drude_1863_1906}
				\\
				{\centering\footnotesize P. Drude}
				\begin{overpic}[width=0.9\textwidth]{Hendrik_Lorentz_1853_1928}
					% Overlay image on top
					\put(0,0){\includegraphics[width=0.4\textwidth]{Alfred_Nobel_1833_1896}}
				\end{overpic}
				\\
				{\centering\footnotesize H. Lorentz }
			\end{figure}
		\end{column}
	\end{columns}
\end{frame}
%
%
\begin{frame}{Independent electron approximation}
	%
\begin{figure}
	\centering
	\includegraphics[width=0.6\textwidth]{metals_independent_electron_approximation}
\end{figure}
	\begin{itemize}
	\item In the condensed phase, there are of the order of \num{e23} electrons.
	\item The repulsive Coulomb potential between an electron and other electron is \textbf{screened} due to large numbers and results in an \textbf{independent} electron.
	\item The problem of motion of \(\mathcal{O}(\num{e23})\) electrons is reduced to the problem of motion of \textbf{single} electron.
\end{itemize}
%
\begin{keyinsight}
	The property of single electron determines the properties of electron gas.
\end{keyinsight}
\end{frame}
\begin{frame}{Free electron approximation}
	%
	\begin{figure}
		\centering
		\includegraphics[width=0.6\textwidth]{metals_free_electron_approximation}
	\end{figure}
	\begin{itemize}
		\item The attractive Coulomb potential between ions and an electron is smeared due to large numbers and results in constant potential in the interior of metal with \textbf{infinite} barrier at the surface. This potential is called \textbf{infinite potential well}.
		\item This potential traps the electrons analogous to gas atoms in a container.
		\item The electrons are free to move in the potential well. The kinetic energy \(K\) of an electron with mass \(m\) and velocity \(\vb{v}\) is
		\[
		K = \frac{1}{2}m v^2
		\]
	\end{itemize}
\end{frame}
%
\begin{frame}{Classical thermodynamics}
	%
\begin{figure}
	\centering
	\includegraphics[width=\textwidth]{metals_classical_thermodynamics}
\end{figure}
\begin{columns}
	\begin{column}[t]{0.85\textwidth}
\begin{footnotesize}
			\begin{itemize}
	\item At temperature \(T=0\), electrons are frozen. At \(T\neq0\), due to the large number of electrons and collisions, the kinetic energy of electrons keeps changing.
	\item However, the probability to find electron with with lower energy is higher than the probality to find electron with higher energy. This is the essence of Maxwell-Boltzmann statistics.
	\item The average kinetic energy \(\langle T \rangle\) at temperature is 
	\[
	\langle K \rangle = \frac{3}{2} k_B T,
	\]
	where \(k_B\) is the Boltzmann constant.
\end{itemize}
\end{footnotesize}
	\end{column}
	%
	\begin{column}[t]{0.13\textwidth}
			\begin{figure}
			\centering
			\includegraphics[width=0.9\textwidth]{Maxwell_Boltzmann_1844_1906}
			\\
			{\centering\footnotesize Boltzmann}
			\includegraphics[width=0.9\textwidth]{James_Clerk_Maxwell_1831_1879}
			\\
			{\centering\footnotesize Maxwell}
		\end{figure}
	\end{column}
\end{columns}
\end{frame}
\begin{frame}{Classical free electron theory}
	%
	\begin{figure}
		\centering
		\href{run:images/electron_drift_lattice_animated_drude_electron_without_electric_field.gif}{%
			\includegraphics[width=0.4\textwidth]{electron_drift_lattice_animated_drude_electron_without_electric_field}%
		}
	\end{figure}
	%
	\begin{itemize}
		\item Electron motion is decomposed into \textbf{thermal motion} and \textbf{drift motion}.
		\item Thermal motion is due to random collisions with the lattice. 
		\item Drift motion is the net motion of electrons. The electron drifts only between collisions.
	\end{itemize}
	%
\end{frame}
%
\begin{frame}{Thermal motion}
	\begin{definition}
		Mean free path \(l_{mf}\) is defined as the average distance the electron travels between collisions.
	\end{definition}
	\begin{definition}
		Thermal velocity \(v_{th}\) is defined as the average velocity of the electron.
	\end{definition}
	%
	\begin{itemize}
		\item Thermal motion has two properties mean free path \(l_{mf}\) and thermal velocity \(v_{th}\).
		\item Drude assumed that \(l_{mf}\) is a material constant.
		\item At room temperature (\SI{~300}{\kelvin}), the thermal velocity is of the \textbf{order of \SI{e5}{\metre\per\second}}.
	\end{itemize}
	%
	\begin{guesstimate}{Average kinetic energy of electron}
		Take room temprerature as \SI{27}{\celsius}. Express energy in \si{\electron\volt}. \hfill [Hint: \SI{1}{\electron\volt} = \SI{1.602e-19}{\joule}]
	\end{guesstimate}
\end{frame}
%
\begin{frame}{Drift motion}
		\begin{itemize}
		\item In the absence of external force, the drift motion is zero.
		\item The external force \(\vb{F}_{ext}\) on the electron is given by the Lorentz force
		\[
		\vb{F}_{ext} = -e \vb{E} -e \vb{v}\times\vb{B},
		\]
		where \(\vb{E}\) is the electric field and \(\vb{B}\) is the magnetic field.
		\item The effect of the collisions is treated as a drag or frictional force
		\[
		\vb{F}_{fric} = -m \frac{\vb{v}}{\tau}.
		\] 
		\item Since Lorentz extended the theory to incorporate electromagnetic force, the theory is also the Lorentz-Drude theory.
	\end{itemize}
\end{frame}
%
\begin{frame}{Lorentz-Drude theory -- Postulates}
	\begin{block}{Postulates}
		\begin{enumerate}
			\item \textbf{Particles:} Electrons are \textbf{classical} particles with mass \(m\) and charge \(-e\).
			\item \textbf{Kinematics:} 
			\begin{itemize}
				\item \textbf{Independent electron approximation:} Electrons are independent and mutual repulsion between them is ignored.

			\end{itemize}
			\item \textbf{Dynamics:} 
			\begin{itemize}
				\item \textbf{Free electron approximation:} Electrons are free and move in an \textbf{infinite potential well}.
				\item \textbf{Thermal motion:} Electron undergoes \textbf{random} collision with ion cores. Random collisions lead to thermal motion.
				\item \textbf{Drift motion:} In between collisions, the Lorentz force acts on electron and leads to drift motion.
			\end{itemize}
			\item \textbf{Classical Thermodynamics:} The average electron energy is given by Maxwell-Boltzmann statistics. 
		\end{enumerate}
	\end{block}
	%
\end{frame}
%
%------------------------
\begin{frame}{Ohm's law}
	\begin{figure}
		\centering
		\href{run:images/electron_drift_lattice_animated_drude_electron_with_electric_field.gif}{%
			\includegraphics[height=0.4\textheight]{electron_drift_lattice_animated_drude_electron_with_electric_field}%
			\includegraphics[height=0.2\textheight]{metals_electron_motion_decomposition}
		}
	\end{figure}
	%
	\begin{itemize}
		\item The motion of electron is decomposed into an \textbf{average} part and the \textbf{random} part.
		\item The average part ``drifts'' with a constant velocity under electric field. The random part has \textbf{zero} net motion.
	\end{itemize}
	\begin{learningobjectives}
		Classical electron theory explains Ohm’s law qualitatively.
	\end{learningobjectives}
\end{frame}
%
\begin{frame}{Ohm's law: Macroscopic version}
	%
	\begin{minipage}[t]{0.7\textwidth}
		\begin{itemize}
			\item The \textbf{macroscopic} version of the Ohm's law states that potential \(V\) applied across a metal is directly proportional to the current \(I\) flowing across it. The constant of proportionaly is called the resistance \(R\).
			\item Its \textbf{microscopic} version relates the electric field and current density
			\item Electric field \(\vb{E}\) is defined as the potential per unit distance. Its SI unit is \si{\volt\per\metre}.
			\item Current density \(\vb{j}\) is defined as the current flowing through a wire per unit area,
			
			where \(A\) is the cross-sectional area of the wire. Its SI unit is \si{\ampere\per\metre\squared}.
		\end{itemize}
	\end{minipage}
	%
	\begin{minipage}[t]{0.28\textwidth}
		\[
		V = IR
		\]
		%\vspace{10ex}
		\begin{figure}
			\centering
			\includegraphics[width=\textwidth]{metals_ohms_law_derivation}
		\end{figure}
		%
		\vspace{3ex}
		\[
		\vb{E} = \frac{V}{\ell}\vu{E}
		\]
		\[
		\vb{j} = \frac{I}{A}\vu{j}
		\]
	\end{minipage}
	%
\end{frame}
%
\begin{frame}{Ohm's law : Microscopic version}
	%
	Lets start with the macroscopic version \(V = R I\) and derive the microscopic version of Ohm's law
	\begin{derivation}
		%
		Upon substitution of \(V, I\) in terms of \(\vb{E}, \vb{j}\), we have
		\[
		\frac{\ell}{V} \vb{E} = \frac{A}{I} \vb{j}
		\quad \Rightarrow \quad
		\vb{j} = \frac{\ell}{A} \frac{I}{V} \vb{E}
		\quad \Rightarrow \quad
		\vb{j} = \frac{\ell}{A R} \vb{E} = \sigma \vb{E}
		\]
		Thus, the microscopic version of Ohm's law is given by
		\begin{center}
			\framebox{$\vb{j} = \sigma \vb{E}$}
		\end{center}
		%
		where \(\sigma\) is given by
		\[
		\sigma = \frac{\ell}{A R}
		\]
	\end{derivation}
\end{frame}
%
\begin{frame}{Conductivity -- \(\sigma\)}
	%
	\begin{definition}
		The conductivity \(\sigma\) is defined as the current density carried by the wire per unit of electric field.
		\[
		\sigma = \frac{j}{E}
		\quad\quad
		\text{SI unit}
		\quad 
		[\sigma] = \left[ \frac{\ell}{A R}\right] = \text{\si{\per\ohm\per\metre}}.
		\]
	\end{definition}
	%
	\begin{question}
		A copper wire of length \(L=\) \SI{2}{\metre} and diameter \(d\)=\SI{1}{\milli\metre} is connected across a DC source. When a voltage of \(V\)=\SI{0.5}{\volt} is applied across the ends of the wire, an electric current of \(I\)=\SI{22}{\ampere} is observed.
		Using these measurements, find the electrical conductivity \(\sigma\) of copper.
	\end{question}
	%
	\begin{keyinsight}
		Metals have high conductivity.
	\end{keyinsight}
\end{frame}
%------------------------
\begin{frame}{Applications of metallic conductivity}
	%
	Due to high conductivity, metals have many applications ranging from power electronics to microelectronics, from mechanical to aerospace industry.
	\begin{figure}
		\centering
		\includegraphics[width=\textwidth]{metals_applications}
	\end{figure}
	\begin{itemize}
		\item Power transmission: Copper/Aluminum cables.
		\item Electronic devices: Contacts, interconnects.
		\item Thermal applications: Heat sinks, cookware.
		\item Alloys: Enhanced strength and corrosion resistance.
	\end{itemize}
\end{frame}

%------------------------
\begin{frame}{Summary of Lecture}
	\begin{itemize}
		\item Materials are classified into conductors, semiconductors, and insulators.
		\item Metals have high conductivity due to free electrons.
		\item Drude's classical electron theory introduced microscopic picture of conduction.
		\item Conductivity is related to resistivity.
	\end{itemize}
\end{frame}
%
\begin{frame}{Materials: Beyond metals}
\begin{tiny}
		\begin{table}
		\centering
		\begin{tabular}{c c c c c c c}
			\toprule
			Condensed	&	Elec.	&	Therm.	&	Opt.	&	Ductile		&	Material	&	Application			\\
			\midrule
			Solid	&	\(\Uparrow\)	&	-	&	-	&	\(\downarrow\)	& Superconductor	&	Levitation\\
			Solid	&	\(\Downarrow\)	&	\(\Uparrow\)	&	-	&	-		& Diamond	&	Chips \\
			Liquid	&	\(\Downarrow\)	&	\(\Uparrow\)	&	-	&	-		& Dielectric coolants	&	Heatsinks \\
			\bottomrule
		\end{tabular}
		\caption{New materials with mixed properties.}
	\end{table}	
\end{tiny}
%
\begin{figure}
	\centering
	\includegraphics[width=\textwidth]{metals_beyond_metals}
\end{figure}
\end{frame}